% !Mode:: "TeX:UTF-8"
% !TEX root = ../../main.tex
% ============================================================
% TOPIC 10
% ============================================================
\mashgulot{10}{Arithmetic operations on vectors: multiplication of a vector by a number, addition and subtraction of vectors}

\begin{tarif}
A quantity determined by its length and its direction is called a
\textbf{vector}. The directed segment with initial point $A$ and terminal
point $B$ is denoted by $\overrightarrow{AB}$. Once a Cartesian coordinate
system is fixed in space, a vector is determined by its coordinates:
$\vec a=(a_1;a_2;a_3)=a_1\vec i+a_2\vec j+a_3\vec k$, where
$\vec i,\vec j,\vec k$ are the unit vectors of the coordinate axes.
\end{tarif}

If $A(x_1;y_1;z_1)$ and $B(x_2;y_2;z_2)$, then
\[\kf{\overrightarrow{AB}=(x_2-x_1;\ y_2-y_1;\ z_2-z_1),\qquad
|\overrightarrow{AB}|=\sqrt{(x_2-x_1)^2+(y_2-y_1)^2+(z_2-z_1)^2}.}\]
The \textbf{length (magnitude)} of $\vec a$ is
$|\vec a|=\sqrt{a_1^2+a_2^2+a_3^2}$.
The vector of length 0 is the \textbf{zero vector} $\vec 0$, and a vector
of length 1 is called a \textbf{unit vector}.

\subsection{Addition, subtraction and multiplication by a number}

\begin{tarif}
For vectors $\vec a=(a_1;a_2;a_3)$, $\vec b=(b_1;b_2;b_3)$ and a number $\lambda$,
\[\kf{\vec a\pm\vec b=(a_1\pm b_1;\ a_2\pm b_2;\ a_3\pm b_3),\qquad
\lambda\vec a=(\lambda a_1;\ \lambda a_2;\ \lambda a_3).}\]
Geometrically, the sum is constructed by the \textbf{triangle rule} or the
\textbf{parallelogram rule}. The vector $\lambda\vec a$ has length
$|\lambda|\,|\vec a|$; it points in the same direction as $\vec a$ when
$\lambda>0$ and in the opposite direction when $\lambda<0$.
\end{tarif}

\textbf{Properties of the linear operations.}
\begin{enumerate}[label=\arabic*)]
\item $\vec a+\vec b=\vec b+\vec a$, \quad $(\vec a+\vec b)+\vec c=\vec a+(\vec b+\vec c)$;
\item $\vec a+\vec 0=\vec a$, \quad $\vec a+(-\vec a)=\vec 0$;
\item $\lambda(\vec a+\vec b)=\lambda\vec a+\lambda\vec b$, \quad $(\lambda+\mu)\vec a=\lambda\vec a+\mu\vec a$;
\item $\lambda(\mu\vec a)=(\lambda\mu)\vec a$, \quad $1\cdot\vec a=\vec a$.
\end{enumerate}

\begin{misol}
Let $\vec a=(2;-1;3)$ and $\vec b=(1;4;-2)$. Find $\vec a+\vec b$, $\vec a-\vec b$,
$3\vec a-2\vec b$ and $|\vec a|$.

\textbf{Solution.}
Adding and subtracting coordinatewise gives
$\vec a+\vec b=(3;3;1)$ and $\vec a-\vec b=(1;-5;5)$. Further,
\[\kf{3\vec a-2\vec b=(6;-3;9)-(2;8;-4)=(4;-11;13),\qquad
|\vec a|=\sqrt{4+1+9}=\sqrt{14}.}\]
\end{misol}

\begin{misol}
Given the points $A(1;2;-1)$ and $B(4;0;5)$, find $\overrightarrow{AB}$, its
length and the midpoint of the segment $AB$.

\textbf{Solution.} We have $\overrightarrow{AB}=(4-1;\ 0-2;\ 5+1)=(3;-2;6)$
and $|\overrightarrow{AB}|=\sqrt{9+4+36}=7$. The midpoint is
$M\left(\frac{1+4}{2};\frac{2+0}{2};\frac{-1+5}{2}\right)=M\left(\frac52;1;2\right)$.
\end{misol}

\subsection{Collinear vectors. Unit vectors}

\begin{tarif}
Vectors lying on one line or on parallel lines are called
\textbf{collinear}. Vectors $\vec a$ and $\vec b\neq\vec 0$ are collinear
if and only if $\vec a=\lambda\vec b$ for some number $\lambda$, that is,
if and only if their coordinates are proportional:
\[\kf{\frac{a_1}{b_1}=\frac{a_2}{b_2}=\frac{a_3}{b_3}.}\]
\end{tarif}

The unit vector in the direction of $\vec a\neq\vec 0$ is
$\vec a^{\,0}=\dfrac{\vec a}{|\vec a|}=(\cos\alpha;\cos\beta;\cos\gamma)$,
where $\alpha,\beta,\gamma$ are the angles that $\vec a$ makes with the
coordinate axes (the direction cosines satisfy
$\cos^2\alpha+\cos^2\beta+\cos^2\gamma=1$).

\begin{misol}
For $\vec a=(2;-6;3)$ we have $|\vec a|=\sqrt{4+36+9}=7$, hence
$\vec a^{\,0}=\left(\frac27;-\frac67;\frac37\right)$.
Since $\vec b=(-4;12;-6)=-2\vec a$, the vectors $\vec a$ and $\vec b$ are
collinear and point in opposite directions.
\end{misol}

\begin{misol}
Find the point $M$ dividing the segment with endpoints $A(1;-2;4)$ and
$B(7;4;-2)$ in the ratio $AM:MB=\lambda=1:2$.

\textbf{Solution.} By the section formula $M=\dfrac{A+\lambda B}{1+\lambda}$
with $\lambda=\frac12$,
$x=\dfrac{1+\frac72}{\frac32}=3$, $y=\dfrac{-2+2}{\frac32}=0$,
$z=\dfrac{4-1}{\frac32}=2$. Thus $M(3;0;2)$.
\end{misol}

\subsection{Applications of vectors in information security}

\textbf{Feature vectors.}
In an intrusion detection system (IDS) each network session is described
by a feature vector $\vec x=(x_1;x_2;\dots;x_n)$ whose coordinates are the
number of packets, the number of bytes transferred, the duration of the
session, the number of errors, and so on. The mean vector of normal
activity, $\vec m=\frac1N(\vec x_1+\dots+\vec x_N)$, is computed with
nothing more than vector addition and multiplication by a number. For a new
session the difference vector $\vec x-\vec m$ is used to detect anomalous
behavior.

\textbf{XOR as addition of binary vectors.}
Binary vectors are added coordinatewise modulo 2 (the XOR operation). The
one-time pad (Vernam cipher) is built on exactly this operation:
$c=p\oplus k$ and $p=c\oplus k$.

\begin{misol}
For the plaintext $p=(1;0;1;1;0)$ and the key $k=(0;1;1;0;1)$ we obtain
$c=p\oplus k=(1;1;0;1;1)$; decryption gives back
$c\oplus k=(1;0;1;1;0)=p$.
\end{misol}

\textbf{Lattices.} In lattice-based cryptography a lattice is the set of all
integer linear combinations of basis vectors $\vec b_1,\dots,\vec b_n$:
$L=\{k_1\vec b_1+\dots+k_n\vec b_n\colon k_i\in\Z\}$. The problem of finding
a shortest nonzero vector of a lattice (SVP) underlies the security of
ciphers designed to resist quantum computers.

\subsection{Computing with Python, Maple and Mathcad}

The following short programs repeat the first two examples of this topic:
linear combinations of vectors, the length of a vector and the midpoint of
a segment.

\begin{Verbatim}[fontsize=\footnotesize,frame=single,label={\small Python}]
import numpy as np
a = np.array([2, -1, 3])
b = np.array([1, 4, -2])
print(a + b, a - b, 3*a - 2*b)
print(np.linalg.norm(a))
A = np.array([1, 2, -1])
B = np.array([4, 0, 5])
print(B - A, np.linalg.norm(B - A), (A + B)/2)
\end{Verbatim}

\begin{Verbatim}[fontsize=\footnotesize,frame=single,label={\small Output}]
[3 3 1] [ 1 -5  5] [  4 -11  13]
3.7416573867739413
[ 3 -2  6] 7.0 [2.5 1.  2. ]
\end{Verbatim}

In Maple, column vectors are entered with angle brackets, and the
Euclidean length is returned by \texttt{Norm(a, 2)}.

\begin{Verbatim}[fontsize=\footnotesize,frame=single,label={\small Maple}]
with(LinearAlgebra):
a := <2, -1, 3>:  b := <1, 4, -2>:
a + b;  a - b;  3*a - 2*b;
Norm(a, 2);                  # sqrt(14)
A := <1, 2, -1>:  B := <4, 0, 5>:
B - A;  Norm(B - A, 2);  (A + B)/2;
\end{Verbatim}

In Mathcad Prime the vectors are typed as one-column matrices; the
absolute-value operator $|\,\cdot\,|$ applied to a vector gives its length.

\begin{Verbatim}[fontsize=\footnotesize,frame=single,label={\small Mathcad}]
a := |  2 |      b := |  1 |
     | -1 |           |  4 |
     |  3 |           | -2 |

3*a - 2*b = |   4 |      |a| = 3.742
            | -11 |
            |  13 |
\end{Verbatim}

\subsection{Exercises on operations with vectors}

\emph{1. Vectors $\vec a$ and $\vec b$ are given. Find the indicated linear combination.}
\bigskip
\begin{enumerate}[label=\arabic*., leftmargin=*]
\item $\vec a=(4;\,-5;\,1)$, $\vec b=(2;\,4;\,-5)$; \ find $-2\vec a+3\vec b$.
\item $\vec a=(2;\,-1;\,5)$, $\vec b=(-3;\,-5;\,3)$; \ find $-\vec a-3\vec b$.
\item $\vec a=(-2;\,0;\,-5)$, $\vec b=(1;\,-3;\,4)$; \ find $-\vec a-\vec b$.
\item $\vec a=(5;\,-1;\,2)$, $\vec b=(-3;\,5;\,-1)$; \ find $2\vec a-\vec b$.
\item $\vec a=(-3;\,2;\,-2)$, $\vec b=(2;\,4;\,1)$; \ find $-3\vec a+\vec b$.
\item $\vec a=(-5;\,-2;\,-3)$, $\vec b=(-2;\,-1;\,3)$; \ find $-\vec a+3\vec b$.
\item $\vec a=(-2;\,0;\,5)$, $\vec b=(3;\,2;\,1)$; \ find $-3\vec a+2\vec b$.
\item $\vec a=(4;\,0;\,3)$, $\vec b=(-3;\,-2;\,1)$; \ find $-2\vec a-3\vec b$.
\item $\vec a=(-5;\,2;\,-1)$, $\vec b=(4;\,5;\,-4)$; \ find $\vec a+3\vec b$.
\item $\vec a=(-4;\,-3;\,1)$, $\vec b=(4;\,0;\,4)$; \ find $-2\vec a-3\vec b$.
\item $\vec a=(-4;\,2;\,-3)$, $\vec b=(-2;\,0;\,1)$; \ find $-2\vec a+2\vec b$.
\item $\vec a=(-1;\,-3;\,4)$, $\vec b=(3;\,-3;\,-4)$; \ find $-\vec a-\vec b$.
\item $\vec a=(-3;\,5;\,5)$, $\vec b=(-3;\,-3;\,2)$; \ find $3\vec a-\vec b$.
\item $\vec a=(0;\,1;\,-2)$, $\vec b=(-5;\,3;\,-5)$; \ find $-\vec a-\vec b$.
\item $\vec a=(-2;\,-4;\,-1)$, $\vec b=(2;\,1;\,4)$; \ find $-2\vec a+3\vec b$.
\item $\vec a=(2;\,5;\,-2)$, $\vec b=(3;\,-1;\,3)$; \ find $\vec a-3\vec b$.
\item $\vec a=(-3;\,2;\,5)$, $\vec b=(2;\,1;\,-3)$; \ find $\vec a+\vec b$.
\item $\vec a=(0;\,3;\,1)$, $\vec b=(2;\,-3;\,3)$; \ find $-3\vec a+3\vec b$.
\item $\vec a=(-5;\,-5;\,2)$, $\vec b=(3;\,-4;\,-5)$; \ find $-\vec a-3\vec b$.
\item $\vec a=(-3;\,-4;\,4)$, $\vec b=(2;\,2;\,-3)$; \ find $3\vec a+3\vec b$.
\item $\vec a=(3;\,2;\,0)$, $\vec b=(-1;\,2;\,1)$; \ find $-\vec a+2\vec b$.
\item $\vec a=(3;\,3;\,-3)$, $\vec b=(-2;\,0;\,4)$; \ find $-3\vec a+2\vec b$.
\item $\vec a=(-2;\,5;\,4)$, $\vec b=(4;\,5;\,2)$; \ find $-3\vec a+3\vec b$.
\item $\vec a=(-4;\,-1;\,1)$, $\vec b=(-2;\,-1;\,-4)$; \ find $-3\vec a+\vec b$.
\item $\vec a=(3;\,-4;\,-1)$, $\vec b=(-3;\,5;\,2)$; \ find $-2\vec a-\vec b$.
\item $\vec a=(4;\,1;\,-2)$, $\vec b=(0;\,-3;\,-4)$; \ find $\vec a-2\vec b$.
\item $\vec a=(1;\,-2;\,-4)$, $\vec b=(5;\,-1;\,2)$; \ find $3\vec a+2\vec b$.
\item $\vec a=(3;\,1;\,2)$, $\vec b=(5;\,3;\,-2)$; \ find $-\vec a+\vec b$.
\item $\vec a=(0;\,-4;\,5)$, $\vec b=(5;\,-5;\,-2)$; \ find $\vec a+3\vec b$.
\item $\vec a=(1;\,5;\,-2)$, $\vec b=(-4;\,1;\,-4)$; \ find $3\vec a-3\vec b$.
\end{enumerate}

\emph{2. For the given points $A$ and $B$, find the vector $\overrightarrow{AB}$, its length and the midpoint $M$ of the segment $AB$.}
\bigskip
\begin{enumerate}[label=\arabic*., leftmargin=*]
\item $A(-1;\,-2;\,-1)$, $B(1;\,0;\,6)$.
\item $A(-3;\,4;\,5)$, $B(-2;\,1;\,5)$.
\item $A(-2;\,2;\,1)$, $B(2;\,-2;\,0)$.
\item $A(1;\,4;\,-1)$, $B(-1;\,5;\,-2)$.
\item $A(1;\,1;\,1)$, $B(3;\,1;\,-5)$.
\item $A(5;\,2;\,5)$, $B(-1;\,6;\,5)$.
\item $A(0;\,6;\,-4)$, $B(3;\,0;\,5)$.
\item $A(3;\,-5;\,3)$, $B(5;\,-2;\,-6)$.
\item $A(3;\,-1;\,4)$, $B(-1;\,-5;\,5)$.
\item $A(0;\,3;\,-1)$, $B(-2;\,3;\,-4)$.
\item $A(-3;\,-4;\,-2)$, $B(-5;\,2;\,6)$.
\item $A(-6;\,-3;\,3)$, $B(5;\,-5;\,0)$.
\item $A(-3;\,3;\,3)$, $B(1;\,6;\,3)$.
\item $A(-6;\,0;\,-3)$, $B(-1;\,3;\,-3)$.
\item $A(2;\,6;\,3)$, $B(6;\,6;\,-1)$.
\item $A(-6;\,4;\,5)$, $B(3;\,-6;\,-6)$.
\item $A(3;\,-1;\,0)$, $B(-1;\,3;\,4)$.
\item $A(-4;\,6;\,-4)$, $B(-5;\,3;\,6)$.
\item $A(-1;\,-4;\,-6)$, $B(5;\,2;\,2)$.
\item $A(6;\,-1;\,-5)$, $B(0;\,1;\,3)$.
\item $A(-1;\,-6;\,-4)$, $B(1;\,5;\,1)$.
\item $A(3;\,1;\,5)$, $B(1;\,3;\,-4)$.
\item $A(-1;\,0;\,-1)$, $B(2;\,2;\,1)$.
\item $A(0;\,-3;\,-5)$, $B(0;\,4;\,-4)$.
\item $A(-6;\,-4;\,4)$, $B(4;\,-3;\,3)$.
\item $A(0;\,6;\,-1)$, $B(2;\,0;\,-6)$.
\item $A(3;\,-6;\,-4)$, $B(-2;\,-3;\,-6)$.
\item $A(3;\,-4;\,-5)$, $B(4;\,3;\,-1)$.
\item $A(-4;\,-2;\,-6)$, $B(0;\,3;\,6)$.
\item $A(0;\,-5;\,-1)$, $B(6;\,-3;\,1)$.
\end{enumerate}

\emph{3. Applied problems.}
\bigskip
\begin{enumerate}[label=\arabic*., leftmargin=*]
\item \textbf{(Unit vector.)} Find the unit vector in the direction of $\vec a=(2;-1;2)$.
\item \textbf{(Collinearity.)} Are the vectors $\vec a=(2;-4;6)$ and $\vec b=(-3;6;-9)$ collinear? How are they directed relative to each other?
\item \textbf{(Parameters.)} For which values of $m$ and $n$ are the vectors $\vec a=(1;m;3)$ and $\vec b=(2;4;n)$ collinear?
\item \textbf{(Parallelogram.)} The points $A(1;0;2)$, $B(3;1;-1)$, $C(4;-2;0)$ are vertices of a parallelogram $ABCD$. Find the vertex $D$.
\item \textbf{(Centroid.)} Find the centroid of the triangle with vertices $A(2;-1;3)$, $B(4;3;-1)$, $C(0;1;1)$.
\item \textbf{(Direction cosines.)} Find the direction cosines of the vector $\vec a=(1;2;-2)$.
\item \textbf{(Feature vectors.)} The login activity of two users over four time slots is given by the vectors $\vec u=(5;2;0;3)$ and $\vec v=(1;4;2;3)$. Find $\vec u-\vec v$ and $\frac12(\vec u+\vec v)$.
\item \textbf{(Decomposition.)} Express the vector $\vec c=(7;-1)$ as a linear combination of $\vec a=(2;1)$ and $\vec b=(1;-1)$.
\item \textbf{(Length.)} Given $|\vec a|=3$, $|\vec b|=4$ and $\vec a\perp\vec b$, find $|\vec a+\vec b|$ and $|\vec a-\vec b|$.
\item \textbf{(Dividing a segment.)} Find the point $M$ that divides the segment with endpoints $A(1;-2;4)$ and $B(7;4;-2)$ in the ratio $AM:MB=1:2$.
\end{enumerate}

